\documentclass[10pt,a4paper]{amsart}
\usepackage{amsmath,amssymb,amsthm,graphicx}
\usepackage{verbatim}
\newenvironment{reference}{\comment}{\endcomment}
%\newenvironment{reference}{}{}
\newenvironment{slogan}{\comment}{\endcomment}
\newenvironment{history}{\comment}{\endcomment}

% For commutative diagrams we use Xy-pic
\usepackage[all]{xy}

% We use 2cell for 2-commutative diagrams.
\xyoption{2cell}
% Theorem environments.
%
\theoremstyle{plain}
\newtheorem{theorem}[subsection]{Theorem}
\newtheorem{proposition}[subsection]{Proposition}
\newtheorem{lemma}[subsection]{Lemma}

\theoremstyle{definition}
\newtheorem{definition}[subsection]{Definition}
\newtheorem{example}[subsection]{Example}
\newtheorem{exercise}[subsection]{Exercise}
\newtheorem{situation}[subsection]{Situation}

\theoremstyle{remark}
\newtheorem{remark}[subsection]{Remark}
\newtheorem{remarks}[subsection]{Remarks}

\numberwithin{equation}{subsection}

% Macros
%
\def\lim{\mathop{\mathrm{lim}}\nolimits}
\def\colim{\mathop{\mathrm{colim}}\nolimits}
\def\Spec{\mathop{\mathrm{Spec}}}
\def\Hom{\mathop{\mathrm{Hom}}\nolimits}
\def\Ext{\mathop{\mathrm{Ext}}\nolimits}
\def\SheafHom{\mathop{\mathcal{H}\!\mathit{om}}\nolimits}
\def\SheafExt{\mathop{\mathcal{E}\!\mathit{xt}}\nolimits}
\def\Sch{\mathit{Sch}}
\def\Mor{\mathop{\mathrm{Mor}}\nolimits}
\def\Ob{\mathop{\mathrm{Ob}}\nolimits}
\def\Sh{\mathop{\mathit{Sh}}\nolimits}
\def\NL{\mathop{N\!L}\nolimits}
\def\CH{\mathop{\mathrm{CH}}\nolimits}
\def\proetale{{pro\text{-}\acute{e}tale}}
\def\etale{{\acute{e}tale}}
\def\QCoh{\mathit{QCoh}}
\def\Ker{\mathop{\mathrm{Ker}}}
\def\Im{\mathop{\mathrm{Im}}}
\def\Coker{\mathop{\mathrm{Coker}}}
\def\Coim{\mathop{\mathrm{Coim}}}

% Boxtimes
%
\DeclareMathSymbol{\boxtimes}{\mathbin}{AMSa}{"02}

%
% Macros for moduli stacks/spaces
%
\def\QCohstack{\mathcal{QC}\!\mathit{oh}}
\def\Cohstack{\mathcal{C}\!\mathit{oh}}
\def\Spacesstack{\mathcal{S}\!\mathit{paces}}
\def\Quotfunctor{\mathrm{Quot}}
\def\Hilbfunctor{\mathrm{Hilb}}
\def\Curvesstack{\mathcal{C}\!\mathit{urves}}
\def\Polarizedstack{\mathcal{P}\!\mathit{olarized}}
\def\Complexesstack{\mathcal{C}\!\mathit{omplexes}}
% \Pic is the operator that assigns to X its picard group, usage \Pic(X)
% \Picardstack_{X/B} denotes the Picard stack of X over B
% \Picardfunctor_{X/B} denotes the Picard functor of X over B
\def\Pic{\mathop{\mathrm{Pic}}\nolimits}
\def\Picardstack{\mathcal{P}\!\mathit{ic}}
\def\Picardfunctor{\mathrm{Pic}}
\def\Deformationcategory{\mathcal{D}\!\mathit{ef}}
\usepackage{cleveref}
\crefname{originalreference}{original reference}{original references}
\Crefname{originalreference}{Original reference}{Original references}
\makeatletter
\newlabel{stacks-sheaves-proposition-exactness-cech-complex}{{06XE}{source}{}{sourcepointer-1}{}}
\newlabel{stacks-sheaves-proposition-exactness-cech-complex@cref}{{[originalreference][0][]06XE}{[1][0][]source}{}{}{}}
\makeatother
\crefname{originalreference}{original reference}{original references}
\Crefname{originalreference}{Original reference}{Original references}
\makeatletter
\newlabel{stacks-sheaves-lemma-surjective-flat-locally-finite-presentation}{{06XH}{source}{}{sourcepointer-2}{}}
\newlabel{stacks-sheaves-lemma-surjective-flat-locally-finite-presentation@cref}{{[originalreference][0][]06XH}{[1][0][]source}{}{}{}}
\makeatother
\crefname{originalreference}{original reference}{original references}
\Crefname{originalreference}{Original reference}{Original references}
\makeatletter
\newlabel{stacks-sheaves-lemma-pullback-lqc}{{06WL}{source}{}{sourcepointer-3}{}}
\newlabel{stacks-sheaves-lemma-pullback-lqc@cref}{{[originalreference][0][]06WL}{[1][0][]source}{}{}{}}
\makeatother
\crefname{originallemma}{original lemma}{original lemmas}
\Crefname{originallemma}{Original lemma}{Original lemmas}
\makeatletter
\newlabel{lemma-pushforward-locally-quasi-coherent}{{075Z}{source}{}{sourcepointer-4}{}}
\newlabel{lemma-pushforward-locally-quasi-coherent@cref}{{[originallemma][0][]075Z}{[1][0][]source}{}{}{}}
\makeatother
\crefname{originalreference}{original reference}{original references}
\Crefname{originalreference}{Original reference}{Original references}
\makeatletter
\newlabel{stacks-sheaves-lemma-lqc-colimits}{{06WM}{source}{}{sourcepointer-5}{}}
\newlabel{stacks-sheaves-lemma-lqc-colimits@cref}{{[originalreference][0][]06WM}{[1][0][]source}{}{}{}}
\makeatother
\crefname{originalreference}{original reference}{original references}
\Crefname{originalreference}{Original reference}{Original references}
\makeatletter
\newlabel{stacks-morphisms-lemma-surjective-family-flat-locally-finite-presentation}{{07AN}{source}{}{sourcepointer-6}{}}
\newlabel{stacks-morphisms-lemma-surjective-family-flat-locally-finite-presentation@cref}{{[originalreference][0][]07AN}{[1][0][]source}{}{}{}}
\makeatother
\title{Local quasi-coherence descends along flat lfp coverings}
\author{The Stacks Project Authors}
\date{stacks-project:stacks-cohomology.tex:local-snapshot-20261012; GNU FDL 1.2 or later}
\begin{document}
\maketitle
\noindent\textit{Editorial note.} All delivered mathematical text is the author's own, copied verbatim from the pinned Stacks Project chapter file stacks-cohomology.tex; no statement or proof step was rewritten, reordered or removed.\par\medskip
\noindent\textit{Editorial note.} The counted span cites the lemmas recorded in assembly.external\_references with canonical tags as navigation pointers.\par\medskip
\noindent\textit{Allowed standard prerequisites (editorial).}\begin{itemize}
\item Exactness of the Cech complex. Sheaves on Stacks, Proposition stacks-sheaves-proposition-exactness-cech-complex with Lemma stacks-sheaves-lemma-surjective-flat-locally-finite-presentation. Provides 0 -> F -> a\_*a\textasciicircum{}*F -> b\_*b\textasciicircum{}*F exact.
\item Pullback preserves lqc. Sheaves on Stacks, Lemma stacks-sheaves-lemma-pullback-lqc. b\textasciicircum{}*F is lqc.
\item Pushforward of lqc is lqc. Cohomology of Algebraic Stacks, Lemma stacks-cohomology-lemma-pushforward-locally-quasi-coherent. a\_*a\textasciicircum{}*F and b\_*b\textasciicircum{}*F are lqc.
\item Surjective flat lfp families. Morphisms of Stacks, Lemma stacks-morphisms-lemma-surjective-family-flat-locally-finite-presentation. Provides the covering factoring through the f\_j.
\end{itemize}\medskip
\ifdefined\hypertarget\hypertarget{sourcepointer-1}{}\fi
\noindent\textit{Source reference.} Original reference 06XE: sources/stacks-chow/tags.txt, source lines 8848--8848. Cross-chapter standard prerequisite cited by the author inside the counted span; it is used as a pointer and its own proof is not part of this item.\par\medskip
\ifdefined\hypertarget\hypertarget{sourcepointer-2}{}\fi
\noindent\textit{Source reference.} Original reference 06XH: sources/stacks-chow/tags.txt, source lines 8851--8851. Cross-chapter standard prerequisite cited by the author inside the counted span; it is used as a pointer and its own proof is not part of this item.\par\medskip
\ifdefined\hypertarget\hypertarget{sourcepointer-3}{}\fi
\noindent\textit{Source reference.} Original reference 06WL: sources/stacks-chow/tags.txt, source lines 8820--8820. Cross-chapter standard prerequisite cited by the author inside the counted span; it is used as a pointer and its own proof is not part of this item.\par\medskip
\ifdefined\hypertarget\hypertarget{sourcepointer-4}{}\fi
\noindent\textit{Source reference.} Original lemma 075Z: sources/stacks-chow/tags.txt, source lines 9148--9148. Cross-chapter standard prerequisite cited by the author inside the counted span; it is used as a pointer and its own proof is not part of this item.\par\medskip
\ifdefined\hypertarget\hypertarget{sourcepointer-5}{}\fi
\noindent\textit{Source reference.} Original reference 06WM: sources/stacks-chow/tags.txt, source lines 8821--8821. Cross-chapter standard prerequisite cited by the author inside the counted span; it is used as a pointer and its own proof is not part of this item.\par\medskip
\ifdefined\hypertarget\hypertarget{sourcepointer-6}{}\fi
\noindent\textit{Source reference.} Original reference 07AN: sources/stacks-chow/tags.txt, source lines 9313--9313. Cross-chapter standard prerequisite cited by the author inside the counted span; it is used as a pointer and its own proof is not part of this item.\par\medskip
\begin{lemma}
\label{lemma-check-lqc-on-flat-covering}
Let $\mathcal{X}$ be an algebraic stack. Let
$f_j : \mathcal{X}_j \to \mathcal{X}$ be a family of flat
and locally finitely presented morphisms of algebraic stacks with
$|\mathcal{X}| =\bigcup |f_j|(|\mathcal{X}_j|)$.
Let $\mathcal{F}$ be a sheaf of $\mathcal{O}_\mathcal{X}$-modules
on $\mathcal{X}_{fppf}$. If each $f_j^{-1}\mathcal{F}$
is locally quasi-coherent, then so is $\mathcal{F}$.
\end{lemma}

\begin{proof}
First, suppose there is a morphism $a : \mathcal{U} \to \mathcal{X}$
which is surjective, flat, locally of finite presentation, quasi-compact,
and quasi-separated such that $a^*\mathcal{F}$ is locally quasi-coherent.
Then there is an exact sequence
$$
0 \to \mathcal{F} \to a_*a^*\mathcal{F} \to b_*b^*\mathcal{F}
$$
where $b$ is the morphism
$b : \mathcal{U} \times_\mathcal{X} \mathcal{U} \to \mathcal{X}$, see
Sheaves on Stacks, Proposition
\ref{stacks-sheaves-proposition-exactness-cech-complex} and
Lemma \ref{stacks-sheaves-lemma-surjective-flat-locally-finite-presentation}.
Moreover, the pullback $b^*\mathcal{F}$ is the pullback of $a^*\mathcal{F}$
via one of the projection morphisms, hence is locally quasi-coherent
(Sheaves on Stacks, Lemma \ref{stacks-sheaves-lemma-pullback-lqc}).
The modules $a_*a^*\mathcal{F}$ and $b_*b^*\mathcal{F}$ are locally
quasi-coherent by Lemma \ref{lemma-pushforward-locally-quasi-coherent}.
(Note that $a_*$ and $b_*$ don't care about which topology is
used to calculate them.)
We conclude that $\mathcal{F}$ is locally quasi-coherent, see
Sheaves on Stacks, Lemma \ref{stacks-sheaves-lemma-lqc-colimits}.

\medskip\noindent
We are going to reduce the proof of the general case the
situation in the first paragraph. Let $x$ be an object of $\mathcal{X}$
lying over the scheme $U$. We have to show that
$\mathcal{F}|_{U_\etale}$ is a quasi-coherent $\mathcal{O}_U$-module.
It suffices to do this (Zariski) locally on $U$, hence we may
assume that $U$ is affine. By
Morphisms of Stacks, Lemma
\ref{stacks-morphisms-lemma-surjective-family-flat-locally-finite-presentation}
there exists an fppf covering $\{a_i : U_i \to U\}$ such that
each $x \circ a_i$ factors through some $f_j$. Hence $a_i^*\mathcal{F}$
is locally quasi-coherent on $(\Sch/U_i)_{fppf}$. After refining
the covering we may assume $\{U_i \to U\}_{i = 1, \ldots, n}$
is a standard fppf covering. Then $x^*\mathcal{F}$ is an fppf
module on $(\Sch/U)_{fppf}$ whose pullback by the morphism
$a : U_1 \amalg \ldots \amalg U_n \to U$ is locally quasi-coherent.
Hence by the first paragraph we see that $x^*\mathcal{F}$ is locally
quasi-coherent, which certainly implies that $\mathcal{F}|_{U_\etale}$
is quasi-coherent.
\end{proof}

\end{document}
